\section{结论}

本笔记考察了将自组织目标围捕控制扩展到异构高阶系统的若干尝试。
主要教训在于：凸包包含部分与避碰部分具有不同级别的难度。
通过利用 APF 项的反对称性，通常可获得
\[
    \sum_i p_{i0}\to0
\]
或该关系的加权变体。
这足以证明目标以渐近或宽松的意义位于智能体的凸包内。
然而，这一中心关系并不蕴含每一对智能体之间的距离始终保持在
安全距离之上。

基于积分的控制器最清晰地演示了这一点。
它可以将每个智能体驱动到 APF 的重复积分，从而将目标置于
智能体的算术中心。
代价是重复积分可能无界增长，因此凸包可能无限扩展。
这仍满足宽松的围捕要求，但它不是有界刚性编队，且本身并不证明避碰。

APF 滤波和指令滤波反步尝试处理了另一个障碍：
对原始 APF 求导对于高阶智能体是不可取的，因为 APF 可以是非光滑的
且在边界处奇异。
滤波使信号可用于高阶控制，但它们也破坏了当 APF 作为力直接进入
物理动力学时可用的直接障碍能量抵消。
因此，APF 相关参考量的有限时间或渐近跟踪并不自动蕴含
安全集合的前向不变性。

增强 APF 函数 $\alpha_r$ 的奇异性本身也不够。
更快的爆破仅当趋近安全边界的速率已被控制时才有帮助。
若无此类速率估计，有限时间趋近可使即使奇异的 APF 信号沿轨迹
变为可积的。
因此，缺失的分量不仅是更强的势函数，而是真正的不变性论证。

未来的工作应聚焦于结合上述尝试中的有用部分。
一个有前景的方向是使用指令滤波或动态面反步来处理异构高阶动力学，
同时添加高阶控制障碍函数或障碍李雅普诺夫层来证明避碰。
这种混合设计将保留自组织 APF 的直觉，同时提供当前尝试所缺失的
安全性证明。